\pdfinfoomitdate=1
\pdftrailerid{}
\documentclass[11pt,letterpaper]{article}
\usepackage[T1]{fontenc}
\usepackage{lmodern}
\usepackage[margin=1in]{geometry}
\usepackage{amsmath,amssymb,amsthm,mathtools,booktabs,microtype}
\usepackage[hidelinks]{hyperref}
\hypersetup{pdftitle={Late growth of Bessel counting coefficients},pdfsubject={A336293 and the rational quotient A395976 over A395977},pdfauthor={Research report}}
\setlength{\emergencystretch}{2em}
\numberwithin{equation}{section}
\newtheorem{theorem}{Theorem}[section]
\newtheorem{lemma}[theorem]{Lemma}
\newtheorem{proposition}[theorem]{Proposition}
\newtheorem{corollary}[theorem]{Corollary}
\theoremstyle{remark}\newtheorem{remark}[theorem]{Remark}
\newcommand{\E}{\mathbb E}
\newcommand{\e}{\mathrm e}
\newcommand{\ii}{\mathrm i}
\newcommand{\Ord}{\mathcal O}
\newcommand{\As}{\mathcal A}
\newcommand{\R}{\mathbb R}
\newcommand{\fall}[2]{(#1)_{\underline{#2}}}
\title{Late growth of Bessel counting coefficients}
\author{Research report}
\date{2 October 2026}
\begin{document}
\maketitle
\begin{abstract}
For the counting sequence $a_n=\sum_{k=0}^n\binom nk^2\binom{2k}k(n-k)!$, let $d_j$ denote the rational coefficients in its normalized expansion in powers of $n^{-1/2}$. We prove the conjectured equivalent $d_j\sim(2\e^2/\pi)\Gamma(j)/4^j$ and give its complete expansion to every fixed inverse-factorial order. Only even Gamma shifts occur. An exact saddle transform connects the coefficients to the square of a classical Bessel asymptotic series. A summable Gaussian lag bound proves the leading constant directly; a stripped-kernel estimate and uniform Cauchy remainder justify all higher orders before factorial-series algebra identifies them. The corrections have a finite Bernoulli-polynomial generator. We also give asymptotic inverses for the original counts and the late coefficients, with integer rounding handled separately. Exact arithmetic and high-precision checks are supplied in a replayable package. The result concerns the quotient A395976/A395977, not separate numerator or denominator growth; it proves eventual positivity but leaves positivity at every index greater than five unresolved.
\end{abstract}

\section{The sequence and the result}
All logarithms are natural. The rising factorial is $(a)_r=a(a+1)\cdots(a+r-1)$, and $\fall{x}{r}=x(x-1)\cdots(x-r+1)$. Empty products equal one. Formal power series need not converge. Every asymptotic remainder is at a fixed finite order, with constants allowed to depend on that order.

The sequence OEIS A336293 is
\begin{equation}\label{eq:an}
 a_n=\sum_{k=0}^n\binom nk^2\binom{2k}k(n-k)!
     =(n!)^2[x^n]\e^x I_0(2\sqrt{x})^2.
\end{equation}
Indeed $I_0(2\sqrt{x})=\sum_{k\ge0}x^k/(k!)^2$, and its square has coefficient $\binom{2k}k/(k!)^2$. Thus the entire generating function in \eqref{eq:an} has no branch ambiguity. The alternative identity
\begin{equation}\label{eq:egf}
 \sum_{n\ge0}\frac{a_n}{n!}z^n
 =\frac1{1-z}\exp\!\left(\frac{2z}{1-z}\right)
     I_0\!\left(\frac{2z}{1-z}\right)
\end{equation}
follows from the binomial transform of $\binom{2k}k/k!$.

The rational coefficients $d_j$ are uniquely specified by
\begin{equation}\label{eq:forward}
 a_n=\frac{n^n\e^{4\sqrt n-n-2}}{\sqrt{8\pi}}
 \left(\sum_{j=0}^{J}d_j n^{-j/2}+\Ord_J(n^{-(J+1)/2})\right)
 \qquad(J\ge0\text{ fixed}).
\end{equation}
Their reduced numerators and positive denominators are recorded in A395976 and A395977. The accessible indexed A395976 entry attributes the late-growth conjecture to Vaclav Kotesovec, 21 May 2026 \cite{oeis336,oeis395976,oeis395977}. In our notation its displayed equivalent is
\begin{equation}\label{eq:oeis}
 d_j\sim\frac{j^{j-1/2}}{\sqrt\pi\,2^{2j-3/2}\e^{j-2}}.
\end{equation}
By Stirling's formula this is exactly the leading term of the following result.

\begin{theorem}\label{thm:main}
The coefficients defined by \eqref{eq:forward} are rational and satisfy, for every fixed integer $M\ge1$,
\begin{equation}\label{eq:main}
 d_j=\frac{2\e^2}{\pi4^j}
 \left\{\sum_{m=0}^{M-1}16^m c_m\Gamma(j-2m)
                  +\Ord_M\bigl(\Gamma(j-2M)\bigr)\right\},\qquad j\to\infty.
\end{equation}
The rational corrections $c_m$ are generated by the finite formula
\begin{equation}\label{eq:cm}
 c_m=\sum_{r=0}^{m}\frac{(1/2)_r^3}{r!4^r}
 [z^{m-r}]\exp\!\left\{
 \sum_{k=1}^{m-r}\frac{(-1)^{k+1}
 [2B_{k+1}(1)-B_{k+1}(r+3/2)]}{k(k+1)}z^k\right\},
\end{equation}
where $B_k(x)$ denotes the Bernoulli polynomial. In particular,
\begin{align}\label{eq:cfirst}
 (c_0,c_1,c_2,c_3,c_4)
 =\left(1,-\frac{25}{96},\frac{2299}{18432},
       -\frac{1073363}{26542080},-\frac{3620999}{2038431744}\right).
\end{align}
Consequently $d_j>0$ for all sufficiently large $j$, and $d_j/d_{j-1}\sim(j-1)/4$.
\end{theorem}
The first visible corrections in Gamma normalization are
\begin{equation}\label{eq:normalized}
 \frac{d_j}{(2\e^2/\pi)4^{-j}\Gamma(j)}
 =1-\frac{25}{6\fall{j-1}{2}}+
 \frac{2299}{72\fall{j-1}{4}}-
 \frac{1073363}{6480\fall{j-1}{6}}+\Ord(j^{-8}).
\end{equation}
There is no shift by one, three, or any other odd integer. Notice that a missing first inverse-factorial correction does not mean that all odd powers of $1/j$ vanish after the falling factorials are expanded.

The theorem gives neither separate asymptotics of the reduced numerator and denominator nor an explicit finite positivity cutoff. In particular it does not settle the entry's stronger question whether $d_j>0$ for every $j>5$. The methods used below are classical Bessel and saddle asymptotics and established factorial-series algebra; no new general theory of resurgence is claimed.

\section{A saddle transform for the actual forward coefficients}
Put $t=n^{-1/2}$ and define
\begin{equation}\label{eq:BH}
 B(t)=\sum_{k\ge0}b_kt^k,\quad b_k=\frac{(1/2)_k^2}{k!4^k},
 \qquad H(t)=B(t)^2=\sum_{k\ge0}h_kt^k.
\end{equation}
Throughout, $u=\ii V$, where $V$ is a standard real Gaussian. An expectation of a formal series means the expectation of each coefficient. Such coefficients are polynomials in $u$, possibly multiplied by $\e^{2u}$, so all the required moments exist. Define
\begin{align}
 R_1(t,u)&=\sum_{m\ge1}\frac{(-1)^m}{m+2}u^{m+2}t^m,\label{eq:R1}\\
 R_2(t,u)&=\sum_{m\ge1}4\binom{1/2}{m+1}u^{m+1}t^m,
 \qquad R=R_1+R_2,\label{eq:R2}\\
 S(t)&=\exp\!\left(2\sum_{r\ge1}
           \frac{B_{2r}}{2r(2r-1)}t^{4r-2}\right).\label{eq:S}
\end{align}
Here $B_{2r}=B_{2r}(0)$ are Bernoulli numbers.

\begin{proposition}\label{prop:transform}
The forward expansion \eqref{eq:forward} exists to every fixed order, and its formal coefficient series is
\begin{equation}\label{eq:transform}
 D(t):=\sum_{j\ge0}d_jt^j=\e^2 S(t)T(H)(t),
\end{equation}
where
\begin{equation}\label{eq:T}
 T(H)(t)=\E\left[\e^{2u}(1+tu)^{-3/2}\e^{R(t,u)}
                       H\!\left(\frac{t}{\sqrt{1+tu}}\right)\right].
\end{equation}
\end{proposition}
\begin{proof}
The standard fixed-order Bessel expansion, uniformly in a closed sector about the positive real axis, gives
\begin{equation}\label{eq:bessel}
 I_0(2\sqrt{x})^2\sim\frac{\e^{4\sqrt{x}}}{4\pi\sqrt{x}}H(x^{-1/2}).
\end{equation}
See DLMF 10.40.1 and its complex-argument error discussion \cite{dlmfbessel}. This use of a fixed-order expansion does not assume anything about the late coefficients to be proved.

To justify extracting the actual coefficients, use Cauchy's formula for the entire function in \eqref{eq:an} on $|x|=n$. The integral representation of $I_0$ implies
\[
 |I_0(2\sqrt{n\e^{\ii\theta}})|\le
 \exp\bigl(2\sqrt n\cos(\theta/2)\bigr),\qquad -\pi\le\theta\le\pi.
\]
Thus the exact integrand has exponential modulus bounded by
$\exp(n\cos\theta+4\sqrt n\cos(\theta/2))$. Outside the arc
$|\theta|\le n^{-2/5}$ there is a loss $c n^{1/5}$ in the real exponent, for some $c>0$. Polynomial normalization cannot offset this loss.

Deform the small arc to the vertical segment $x=n+\ii\sqrt n\,v$, joined horizontally to the old endpoints. The deformation lies in an analytic region away from zero. On a joining segment, write $x=n+z$. Then $|\Im z|\asymp n^{3/5}$ and $|\Re z|=\Ord(n^{1/5})$, and
\[
 \Re(x-n\log x)=n-n\log n+\frac{\Re(z^2)}{2n}
                            +\Ord(|z|^3/n^2).
\]
The quadratic term is $-\tfrac12n^{1/5}+o(n^{1/5})$, while the cubic error is $\Ord(n^{-1/5})$ and $4\Re\sqrt{x}-4\sqrt n=o(1)$. The same exact Bessel bound with $\Re\sqrt{x}$ controls the joining pieces. They are exponentially small.

On the vertical segment $|v|=\Ord(n^{1/10})$, put $u=\ii v$. Direct Taylor expansion yields
\begin{align*}
 x-n\log x&=n-n\log n+u^2/2+R_1(t,u),\\
 4\sqrt{x}&=4/t+2u+R_2(t,u).
\end{align*}
Here $|tu|=\Ord(n^{-2/5})$ and $t|v|^3=\Ord(n^{-1/5})$. At each fixed degree, the logarithm, square-root, and amplitude remainders are bounded by the next power of $t$ times a fixed polynomial in $|v|$, times a bounded exponential. The factor $\exp(-v^2/2)$ makes these bounds integrable. The local Bessel remainder is uniform, and polynomial Gaussian tails permit extension to the whole vertical Gaussian line coefficient by coefficient.

The factors $dx/x$ and $x^{-1/2}$ in \eqref{eq:bessel} give $(1+tu)^{-3/2}$. Stirling's expansion of $(n!)^2$ gives $S(t)$. Combining constants gives $\sqrt{2\pi}/(4\pi)=1/\sqrt{8\pi}$, and $\E\e^{2\ii V}=\e^{-2}$ accounts for the exponential $\e^{-2}$ in \eqref{eq:forward}. This proves \eqref{eq:transform}, including its normalization.
\end{proof}

\subsection{An exact coefficient algorithm}
The expectation in \eqref{eq:T} is also an effective rational algorithm. Define tilted moments
\[
 \mu_p=\e^2\E(\e^{2u}u^p),\qquad
 \mu_0=1,\quad\mu_1=-2,\quad
 \mu_p=-2\mu_{p-1}-(p-1)\mu_{p-2}\quad(p\ge2).
\]
They follow by differentiating the Gaussian characteristic function. Expand the kernel to degree $J$, replace each $u^p$ by $\mu_p$, and multiply by the truncation of $S$. All steps are rational and finite. For example,
\begin{align}\label{eq:dsmall}
 (d_0,\ldots,d_6)=\bigg(&1,\frac{55}{24},\frac{2581}{1152},
 \frac{255893}{414720},-\frac{4579919}{7962624},\nonumber\\
 &-\frac{1075995491}{6688604160},
 \frac{2191041054791}{4815794995200}\bigg).
\end{align}
The two negative entries underscore why eventual positivity must be stated carefully.

\section{An elementary proof of the leading late constant}
This section proves the leading conjecture without relying on a composition theorem for factorial series.

\subsection{The input and its convolution}
The exact formula
\begin{equation}\label{eq:bexact}
 b_j=\frac{\Gamma(j+1/2)^2}{\pi\Gamma(j+1)4^j}
\end{equation}
implies $b_j\sim\Gamma(j)/(\pi4^j)$ and uniform two-sided bounds on the ratio for all integer $j\ge1$. The endpoints in $h_j=\sum_{k=0}^j b_kb_{j-k}$ contribute $2b_j$. The interior sum is bounded using
\[
 \sum_{k=1}^{j-1}\frac{\Gamma(k)\Gamma(j-k)}{\Gamma(j)}
 =\frac1{j-1}\sum_{r=0}^{j-2}\binom{j-2}{r}^{-1}=\Ord(j^{-1}).
\]
For $j\ge3$, the last reciprocal-binomial sum is at most three: the endpoints give two and each interior term is at most $1/(j-2)$. Therefore
\begin{equation}\label{eq:hasym}
 h_j\sim\frac2\pi4^{-j}\Gamma(j),\qquad
 C_1 4^{-j}\Gamma(j)\le h_j\le C_2 4^{-j}\Gamma(j)\quad(j\ge1),
\end{equation}
with fixed positive constants $C_1,C_2$.

\subsection{A summable bound over every lag}
Write $t_j=[t^j]T(H)$ and $k=j-\ell\ge1$. The contribution of $h_k$ is
\begin{equation}\label{eq:lag}
 h_k\E\left[\e^{2u}[t^\ell](1+tu)^{-(k+3)/2}\e^{R_1}\e^{R_2}\right].
\end{equation}
Split the degree as $a+b+c=\ell$, with $r$ phase factors from $R_1$ and $s$ from $R_2$. The amplitude coefficient is
\[
 (-1)^a\frac{((k+3)/2)_a}{a!}u^a.
\]
For positive $b$, the absolute sum of coefficients supplied by $r$ factors of $R_1$ is at most
$\binom{b-1}{r-1}|u|^{b+2r}/r!$. For positive $c$, the analogous bound is
$4^s\binom{c-1}{s-1}|u|^{c+s}/s!$. The zero-degree cases each have one term and zero phase factors. These binomial counts are the counts of ordered positive compositions; no uncontrolled factorial count is hidden.

The total moment degree $p=\ell+2r+s$ is at most $3\ell$. Since $|\e^{2u}|=1$, Gaussian absolute moments give
$\E|V|^p\le C^\ell\ell^{p/2}$. All composition counts, powers of four and the polynomial number of tuples can be absorbed in $C^\ell$. The following uniform bounds hold:
\begin{equation}\label{eq:uniformratios}
 \frac{\Gamma(j-\ell)}{\Gamma(j)}\le\e^{2\ell}j^{-\ell},\qquad
 ((k+3)/2)_a\le(Cj)^a,\qquad
 \frac1{a!r!s!}\le\e^{3\ell}\ell^{-a-r-s}.
\end{equation}
For the first inequality, use factors at least $j/2$ when $\ell\le j/2$. Otherwise
\[
 \sum_{q=j-\ell}^{j-1}\log(j/q)
 \le\int_0^j\log(j/x)\,dx=j\le2\ell.
\]
The final inequality in \eqref{eq:uniformratios} follows from $\ell^m/m!\le\e^\ell$ for each of $m=a,r,s$.

Combining these estimates and \eqref{eq:hasym}, the absolute contribution at lag $\ell$, normalized by $h_j$, is bounded by
\begin{align}\label{eq:majorant}
 C^\ell j^{a-\ell}\ell^{\ell/2-a-s/2}
 &=C^\ell\ell^{-\ell/2}(\ell/j)^{b+c}\ell^{-s/2}
 \le M_\ell:=C^\ell\ell^{-\ell/2}.
\end{align}
Take $M_0$ separately. The sequence $(M_\ell)$ and each fixed polynomial multiple are summable. The estimate covers small positive $k$ as well as fixed lag.

The endpoint $k=0$ must be separate because $\Gamma(0)$ is undefined. Repeating the same coefficient bounds with $\ell=j$ and using
$1/h_j\le C^j j^{1-j}$ gives at most $jC^j j^{-j/2}$ after normalization. It is smaller than every power of $1/j$.

\subsection{Fixed lag and the Stirling factor}
For each fixed $\ell$, only $a=\ell$, $b=c=0$ has the full degree $\ell$ in $j$. Using $h_{j-\ell}/h_j\sim4^\ell j^{-\ell}$, its limit is
\[
 \E\left[\e^{2u}\frac{(-2u)^\ell}{\ell!}\right].
\]
The summable bound \eqref{eq:majorant} permits dominated convergence over the triangular array of lags. Consequently
\begin{equation}\label{eq:tlimit}
 \frac{t_j}{h_j}\longrightarrow
 \sum_{\ell\ge0}\E\left[\e^{2u}\frac{(-2u)^\ell}{\ell!}\right]
 =\E(\e^{2u}\e^{-2u})=1.
\end{equation}
The last interchange uses $\E\exp(2|V|)<\infty$.

If $s_p=[t^p]S(t)$, the Bernoulli bound
$|B_{2r}|\le C(2r)!/(2\pi)^{2r}$ gives
\begin{equation}\label{eq:Sbound}
 |s_p|\le C A^p\Gamma(p/2+1).
\end{equation}
For completeness, coefficients of $\log S$ are bounded by an exponential times $\sqrt{p!}$. In a product with total degree $p$, the product of the square roots of the component factorials is at most $\sqrt{p!}$ by multinomial integrality. The number of ordered compositions is at most exponential. Exponentiation preserves this bound, and Stirling's formula converts $\sqrt{p!}$ to an exponential times $\Gamma(p/2+1)$.

The lag bound also gives $|t_j|\le C h_j$ for $j\ge1$. For $1\le p\le j-1$, the normalized contribution of $s_pt_{j-p}$ is at most
\[
 C^p\Gamma(p/2+1)\frac{\Gamma(j-p)}{\Gamma(j)}
 \le C^p p^{p/2}j^{-p}\le C^p p^{-p/2}.
\]
Each fixed positive $p$ vanishes in the limit, and the last majorant is summable. The endpoint $p=j$ is negligible by \eqref{eq:Sbound}. Hence $[t^j]ST/h_j\to1$. Together with \eqref{eq:transform} and \eqref{eq:hasym}, this proves
\begin{equation}\label{eq:leading}
 d_j\sim\frac{2\e^2}{\pi}\frac{\Gamma(j)}{4^j}.
\end{equation}

\section{Factorial expansions with an actual input remainder}
For this proof, say that $f(t)=\sum f_jt^j$ belongs to the factorial class if, for every fixed $R\ge0$,
\begin{equation}\label{eq:class}
 f_j=\sum_{r=0}^{R-1}\alpha_r4^{r-j}\Gamma(j-r)
           +\Ord_R(4^{R-j}\Gamma(j-R)).
\end{equation}
Its asymptotic derivation is $\As f(t)=\sum_{r\ge0}\alpha_rt^r$. This is Borinsky's convention with parameters $(\alpha,\beta)=(1/4,0)$ \cite{borinsky}. Finite small indices are immaterial.

A direct integral verifies membership for $B$, including remainders. Let
$G_n=\Gamma(n+1/2)^2/\Gamma(n+1)$. The beta integral gives
\begin{equation}\label{eq:beta}
 G_n=\frac1{\sqrt\pi}\int_0^\infty v^{n-1/2}\e^{-v}
       \int_0^\infty u^{-1/2}(1+u)^{-1/2}\e^{-vu}\,du\,dv.
\end{equation}
Integrating $v$ first verifies the identity using $B(1/2,n+1/2)$. The degree $R-1$ Taylor polynomial of $(1+u)^{-1/2}$ has absolute remainder bounded on the \emph{entire} half-line $u\ge0$ by $(1/2)_R u^R/R!$. This follows from its derivative bound and the integral Taylor remainder. Integrating its monomials and remainder in \eqref{eq:beta}, for $n>R$, gives
\begin{equation}\label{eq:betarem}
 \left|G_n-\sum_{r=0}^{R-1}\frac{(-1)^r(1/2)_r^2}{r!}\Gamma(n-r)\right|
 \le\frac{(1/2)_R^2}{R!}\Gamma(n-R).
\end{equation}
There is thus no unproved local expansion under an unbounded integral. Dividing by $\pi4^n$ proves
\begin{equation}\label{eq:AB}
 \As B(t)=\frac1\pi B(-t).
\end{equation}

Borinsky's product rule, Proposition 22, and composition theorem, Theorem 35, apply for real $\beta$, including zero; Lemma 37 explicitly establishes the extension from positive $\beta$. In the present normalization they give
\begin{align}\label{eq:algebra}
 \As(fg)&=f\As g+g\As f,\\
 \As(f\circ g)&=(f'\circ g)\As g+
       \exp(4/t-4/g(t))(\As f)(g(t)),\label{eq:chain}
\end{align}
when $g(0)=0$ and $g'(0)=1$. Analytic germs have zero asymptotic derivation. In particular
\begin{equation}\label{eq:AH}
 \As H(t)=\frac2\pi B(t)B(-t).
\end{equation}
These rules identify the eventual answer, but do not by themselves justify an expectation over an unbounded Gaussian parameter. We establish that missing uniformity next.

\section{Uniform transfer to every factorial order}
\subsection{A bound that survives stripping rational factors}
Let $E_d(u)=[t^d]\exp R(t,u)$ and write $a=\ell-d$. Since
$|[t^m]R|\le|u|^{m+2}+|u|^{m+1}$, for $d\ge1$,
\begin{equation}\label{eq:Ed}
 |E_d(u)|\le\sum_{f=1}^d\binom{d-1}{f-1}\frac{2^f}{f!}
                         \max(1,|u|)^{d+2f},\qquad E_0=1.
\end{equation}
For $0\le f\le\ell$, Gaussian moments imply
\begin{equation}\label{eq:gaussmoment}
 \frac{\E\max(1,|V|)^{\ell+2f}}{f!}
 \le C8^\ell\Gamma((\ell+1)/2).
\end{equation}
Indeed $\E|V|^p=2^{p/2}\Gamma((p+1)/2)/\sqrt\pi$, and, with $z=(\ell+1)/2$,
$\Gamma(z+f)/(\Gamma(z)f!)\le2^{z+f}$. The latter inequality follows by bounding one positive term of $(1-1/2)^{-z}$ by its sum. The replacement of $|V|$ by $\max(1,|V|)$ adds at most one.

Using $1/a!\le(\ell+1)^d/\ell!$, and absorbing composition counts into an exponential, yields the \emph{stripped} bound
\begin{equation}\label{eq:stripped}
 \frac{4^\ell}{a!}\E\bigl[|u|^a|E_d(u)|\bigr]
 \le N_\ell(\ell+1)^d,
 \qquad
 N_\ell=C A^\ell(\ell+2)^2\frac{\Gamma((\ell+1)/2)}{\ell!}.
\end{equation}
The constants are independent of $j,\ell,d$. Every polynomial multiple of $N_\ell$ is summable, and its tail beyond $j^\delta$ is smaller than any power of $j^{-1}$ for each fixed $\delta>0$. This follows from $\log N_\ell=-\tfrac12\ell\log\ell+\Ord(\ell)$. Crucially, \eqref{eq:stripped} does not assume that the rational factors about to be expanded are still present.

\subsection{The finite rational product and its remainder}
Fix a target order $R$. Insert one factorial input shift $q<R$ from \eqref{eq:class} for $h_k$, where $k=j-\ell$. After normalizing by $4^{-j}\Gamma(j)$, the exact lag and degree term is a constant depending only on $q$ times
\begin{equation}\label{eq:normalizedkernel}
 \frac{4^\ell(-1)^a2^{-a}}{a!}
 \E\bigl[\e^{2u}u^aE_d(u)\bigr]
 j^{-d-q}P_{\ell,a,q}(1/j),
\end{equation}
where
\begin{equation}\label{eq:P}
 P_{\ell,a,q}(z)=
 \frac{\prod_{v=0}^{a-1}(1+(2v-\ell+3)z)}
      {\prod_{v=1}^{\ell+q}(1-vz)}.
\end{equation}
The numerator comes from removing $(j/2)^a$ from $((k+3)/2)_a$; the denominator comes from removing $j^{-\ell-q}$ from $\Gamma(k-q)/\Gamma(j)$. These give exactly the net power $j^{-d-q}$.

Set $L=\ell+R+3$ and $\rho=(100L^2)^{-1}$. There are at most $2L$ linear factors in \eqref{eq:P}, each with a coefficient of modulus at most $3L$. On $|z|\le\rho$, use the analytic logarithm at one and $|\log(1+w)|\le2|w|$ to obtain
\[
 |\log P(z)|\le12L^2|z|\le12/100,
 \qquad |P(z)|\le\e.
\]
Cauchy's estimate gives $|[z^m]P|\le\e(100L^2)^m$. For $\ell\le j^\delta$ with fixed $0<\delta<1/2$, eventually $j^{-1}\le\rho/2$, uniformly in $\ell$. Therefore
\begin{equation}\label{eq:cauchy}
 \left|P(1/j)-\sum_{m=0}^{K}[z^m]P\,j^{-m}\right|
 \le2\e(100L^2)^{K+1}j^{-K-1}.
\end{equation}
This is a uniform finite-order estimate, not a pointwise expansion with a lag-dependent uncontrolled error.

\subsection{Summing all errors}
Split the original lag sum at $j^\delta$. Its large-lag tail is superpolynomially small by \eqref{eq:majorant}; the fixed small-input endpoints, including $k=0$, satisfy the same conclusion separately. On the retained range $k\ge j/2$ eventually. The remainder of the factorial expansion of $h_k$, available from \eqref{eq:AH} and the product theorem, is controlled by the original absolute kernel bound because
\[
 \Gamma(k-R)/\Gamma(k)=\Ord_R(j^{-R}).
\]
Its total is $\Ord_R(j^{-R})$ after normalization.

For $d+q\ge R$, \eqref{eq:stripped} and $|P(1/j)|\le\e$ give the factor
\[
 j^{-d-q}(\ell+1)^d
 =j^{-R}(\ell+1)^{R-q}
      \left(\frac{\ell+1}{j}\right)^{d+q-R}
 \le j^{-R}(\ell+1)^R.
\]
There are at most $\Ord_R(\ell+1)$ pairs $(d,q)$, so their total error at lag $\ell$ is bounded by
$C_Rj^{-R}N_\ell L^{R+1}$, a summable bound.

For $d+q<R$, choose $K=R-1-d-q$ in \eqref{eq:cauchy}. The outside power $j^{-d-q}$ converts its remainder to exactly $j^{-R}$. The remaining polynomial loss is at most $L^{2R}$ because
$d+2K+2=2R-d-2q\le2R$. There are only boundedly many retained $(d,q)$ pairs, depending on $R$. Hence their summed remainder is also $\Ord_R(j^{-R})$. Each extracted coefficient has a polynomially weighted $N_\ell$ envelope, so its infinite lag sum and Gaussian expectation can be interchanged absolutely.

We have proved an expansion to every order in $1/j$ after normalization by $4^{-j}\Gamma(j)$. The finite triangular relation
\[
 \frac{\Gamma(j-r)}{\Gamma(j)}=\frac1{(j-1)\cdots(j-r)}
\]
converts it, at each finite order, to \eqref{eq:class}. Thus $T(H)$ belongs to the required factorial class \emph{before} the algebra is used to identify its coefficients.

The Stirling factor also belongs to that class, with $\As S=0$: \eqref{eq:Sbound} implies, for every fixed $R$,
\[
 s_p=o\bigl(4^{R-p}\Gamma(p-R)\bigr),
\]
as the logarithm of their absolute ratio is at most $-\tfrac12p\log p+\Ord_R(p)$. The product rule therefore applies to \eqref{eq:transform} without a further interchange: $\As(S T(H))=S\As T(H)$.

\section{Identifying the even corrections}
For fixed real $u$, the multiplier
$K_u(t)=\e^{2u}(1+tu)^{-3/2}\e^{R(t,u)}$ and substitution
$g_u(t)=t/\sqrt{1+tu}$ are analytic germs. The substitution is tangent to the identity. Equations \eqref{eq:chain} and \eqref{eq:AH} identify the extracted factorial expansion of $K_uH(g_u)$. At a fixed asymptotic order, its lag sums are entire in $u$: only bounded phase and shift degrees occur, while the remaining terms have factorial denominators $(\ell-d)!$ and polynomial losses in $\ell$. Their convergence is locally uniform in $u$. The identity theorem extends the formula from real $u$ to imaginary $u$, and the bounds in Section 5 justify Gaussian integration of these coefficient sums.

The exponential correction cancels exactly:
\[
 \frac4t-\frac4{g_u(t)}=-\frac{4(\sqrt{1+tu}-1)}t=-2u-R_2(t,u).
\]
Consequently
\begin{equation}\label{eq:AD}
 \As D(t)=\frac{2\e^2}{\pi}S(t)
 \E\left[(1+tu)^{-3/2}\e^{R_1(t,u)}
                 B(g_u(t))B(-g_u(t))\right].
\end{equation}
This is now a justified identity of formal series, rather than a formal Stokes prescription passed through an unchecked integral.

The classical mixed Bessel product gives
\begin{equation}\label{eq:mixed}
 B(t)B(-t)=\sum_{r\ge0}\beta_rt^{2r},
 \qquad \beta_r=\frac{(1/2)_r^3}{r!4^r}.
\end{equation}
It can be verified either from the product $I_0K_0$ in DLMF 10.40.6 or by its symmetric-square differential equation \cite{dlmfbessel}. Equation \eqref{eq:AD} is invariant under $(t,u)\mapsto(-t,-u)$ after Gaussian expectation, and $S$ is even. Its odd coefficients therefore vanish. Define its even coefficients by
\begin{equation}\label{eq:cgaussian}
 \sum_{m\ge0}c_mt^{2m}
 =S(t)\E\left[(1+tu)^{-3/2}\e^{R_1(t,u)}
          \sum_{r\ge0}\beta_r\frac{t^{2r}}{(1+tu)^r}\right].
\end{equation}
At every fixed degree only finitely many $r$ occur. The established factorial-class remainder and \eqref{eq:AD} now prove \eqref{eq:main}. It remains to give a convenient finite generator equivalent to \eqref{eq:cgaussian}.

\subsection{The reciprocal Gamma interpretation}
For each fixed $r\ge0$, put $s=n+r+3/2$. The reciprocal-Gamma Hankel integral \cite{dlmfgamma} deforms to the upward line $\Re x=n$:
\begin{equation}\label{eq:hankel}
 \frac1{\Gamma(s)}=\frac1{2\pi\ii}
 \int_{n-\ii\infty}^{n+\ii\infty}\e^x x^{-s}\,dx,\qquad s>1.
\end{equation}
Use the branch of $\log x$ real on the positive axis. The integral is absolutely convergent; closing to the left justifies deformation from the standard Hankel loop, with the exponential and the power controlling distant connectors. This is not a Taylor coefficient of $\e^x x^{-r-1/2}$ at the origin, which would be invalid.

On substituting $x=n+\ii\sqrt n\,v$, the exact phase remainder is
\begin{equation}\label{eq:Rexact}
 R_{1,\mathrm{ex}}(t,u)
 =t^{-2}\{tu-\log(1+tu)\}-u^2/2.
\end{equation}
Its local power series is \eqref{eq:R1}, but that series is not asserted to converge for all Gaussian values of $u$. The exact integral becomes
\begin{equation}\label{eq:gammaexact}
 \frac1{\Gamma(n+r+3/2)}=\frac{\e^n n^{-n-r-1}}{\sqrt{2\pi}}
 \E\left[(1+tu)^{-r-3/2}\e^{R_{1,\mathrm{ex}}(t,u)}\right].
\end{equation}
Relative to $\e^n n^{-s}$, the exact vertical integrand's modulus is
$(1+v^2/n)^{-s/2}$. For $|v|\le\sqrt n$ this is at most $\e^{-v^2/4}$. For $|v|>\sqrt n$, set $v=\sqrt n z$ and use $1+z^2\ge2|z|$ to bound the tail by a polynomial factor times $2^{-s/2}$. Thus restriction to $|v|\le n^{1/10}$ loses less than every algebraic order. The local Taylor argument of Section 2 proves the coefficientwise asymptotic identity with $R_1$ in place of $R_{1,\mathrm{ex}}$.

Multiplication by $(n!)^2/(\sqrt{2\pi}n^n\e^{-n})$ gives the formal interpretation
\begin{equation}\label{eq:gammagenerator}
 \sum_{m\ge0}c_m n^{-m}\sim
 \frac{(n!)^2}{\sqrt{2\pi}n^n\e^{-n}}
 \sum_{r\ge0}\frac{\beta_r}{\Gamma(n+r+3/2)}.
\end{equation}
The infinite sum on the right is a coefficient generator: to determine the coefficient of $n^{-m}$ retain only $0\le r\le m$. No convergence or equality of that infinite $r$ sum is claimed.

The logarithmic Gamma expansion with Bernoulli polynomials \cite{dlmfstirling} shows that the fixed-$r$ summand, including its normalization, is
\[
 n^{-r}\exp\!\left\{\sum_{k\ge1}
 \frac{(-1)^{k+1}[2B_{k+1}(1)-B_{k+1}(r+3/2)]}{k(k+1)n^k}\right\}
\]
as a fixed-order expansion. This proves the finite generator \eqref{eq:cm}, and hence \eqref{eq:cfirst}.

\section{Asymptotic inverses and integer thresholds}
An inverse of a smooth asymptotic model and an integer threshold answer different questions. A subunit approximation to the former does not remove rounding uncertainty in the latter. All formulas below use the positive real branch of Lambert $W$ for large positive arguments.

\subsection{Inverting the original counts}
The first forward terms imply
\begin{equation}\label{eq:loga}
 \log a_n=n(\log n-1)+4\sqrt n-2-\tfrac12\log(8\pi)
             +\frac{55}{24\sqrt n}-\frac{37}{96n}+\Ord(n^{-3/2}).
\end{equation}
For a target $x\to\infty$, set
\begin{equation}\label{eq:inverseN}
 Y=\log x+2+\tfrac12\log(8\pi),\qquad
 N=\frac{Y}{W(Y/\e)},\qquad L=\log N.
\end{equation}
Then $N(\log N-1)=Y$. Define
\begin{align}\label{eq:ra}
 r_A(x)={}&N-\frac{4\sqrt N}{L}+\frac8{L^2}-\frac8{L^3}\nonumber\\
 &+\frac1{\sqrt N}\left(-\frac{55}{24L}-\frac8{L^3}
                       +\frac{112}{3L^4}-\frac{32}{L^5}\right).
\end{align}
The inverse of the displayed finite logarithmic model in \eqref{eq:loga}, without its remainder, equals
\begin{equation}\label{eq:raerror}
 r_A(x)+\Ord\bigl((N\log N)^{-1}\bigr).
\end{equation}
The model is continued to sufficiently large real $n$, where it is strictly increasing. The remainder in \eqref{eq:loga} is used separately to compare it with the discrete counts.

To verify the formula, put $n=N+A\sqrt N+B+C/\sqrt N$. The coefficients in the equation through order $N^{-1/2}$ give
\begin{align*}
 A&=-4/L,\qquad B=8/L^2-8/L^3,\\
 C&=-\{AB-A^3/6+2B-A^2/2+55/24\}/L.
\end{align*}
This is the coefficient in \eqref{eq:ra}. The remaining logarithmic residual is $\Ord(N^{-1})$, while the derivative of the finite model is asymptotic to $L$. The mean value theorem proves \eqref{eq:raerror}. Further terms follow recursively by Taylor expansion about $N$, using additional $d_j$ from the exact algorithm.

The counts are strictly increasing. Indeed
$a_n/n!=\sum_{k=0}^n\binom nk\binom{2k}k/k!$ is nondecreasing, so $a_n\ge n a_{n-1}$ for $n\ge2$, and $a_1=3>a_0=1$. Let
\[
 N_A(x)=\min\{n\ge0:a_n\ge x\}.
\]
There exist $C>0$ and $x_0$ such that for $x\ge x_0$,
\begin{equation}\label{eq:thresholda}
 \left\lceil r_A(x)-\frac C{N\log N}\right\rceil
 \le N_A(x)\le
 \left\lceil r_A(x)+\frac C{N\log N}\right\rceil.
\end{equation}
For a proof, apply the uniform logarithmic remainder at the nearby integer indices, compare with the increasing smooth model, and enlarge $C$ to absorb the conversion by its derivative. Indices below the lower real bound have count less than $x$, and an integer at or above the upper bound has count at least $x$. Equation \eqref{eq:thresholda} implies $N_A(x)=r_A(x)+\Ord(1)$. It is an asymptotic enclosure with unspecified constants, not a certified finite-target rounding rule.

\subsection{Inverting the late coefficients}
Theorem \ref{thm:main} and Stirling's formula give
\begin{equation}\label{eq:logd}
 \log d_j=j(\log j-1-\log4)-\tfrac12\log j+K+
                  \frac1{12j}+\Ord(j^{-2}),
 \quad K=\log\!\left(\frac{2\e^2\sqrt{2\pi}}\pi\right).
\end{equation}
Let
\begin{equation}\label{eq:J}
 Z=\log x-K,\qquad J=\frac Z{W(Z/(4\e))},\qquad
 L_d=\log(J/4),\qquad
 r_d(x)=J+\frac{\log J}{2L_d}.
\end{equation}
Then the corresponding smooth inverse is $r_d(x)+o(1)$. More specifically, it is $r_d(x)+\Ord((J L_d)^{-1})$. To see the stated rate, put $\delta=\log J/(2L_d)=\Ord(1)$. The residual after substitution into \eqref{eq:logd} is
\[
 \frac{\delta^2/2-\delta/2+1/12}{J}+\Ord(J^{-2}),
\]
and the derivative is asymptotic to $L_d$. This also gives the optional next Newton term by subtracting the displayed numerator divided by $J L_d$. Arbitrarily many smooth-model corrections can be generated from \eqref{eq:main} and the logarithmic Gamma series; no analytic interpolation of the discrete coefficients is required.

Choose any fixed $j_0$ after which $d_j$ is positive and increasing, whose existence follows from Theorem \ref{thm:main}. The eventual threshold
$N_d(x)=\min\{j\ge j_0:d_j\ge x\}$ is well defined for large $x$. For some $C'>0$ and all sufficiently large $x$,
\begin{equation}\label{eq:thresholdd}
 \left\lceil r_d(x)-\frac{C'}{J L_d}\right\rceil
 \le N_d(x)\le
 \left\lceil r_d(x)+\frac{C'}{J L_d}\right\rceil.
\end{equation}
The same comparison argument proves this enclosure. Hence $N_d(x)=r_d(x)+\Ord(1)$, rather than an unqualified $o(1)$ claim for an integer-valued inverse. For large targets a finite initial set of coefficients has no effect. Neither $j_0$ nor the enclosure constants are made effective by this proof.

\section{A fixed number of colors}
The same leading argument has a useful counting generalization. For each fixed positive integer $q$, define
\[
 A_{q,n}=(n!)^2[x^n]\e^x I_0(2\sqrt{x})^q.
\]
Take two distinguishable labeled $n$-element sets. Partition each into labeled classes $0,1,\ldots,q$ with matching class sizes, and place a bijection between the class-zero subsets. For a size vector $(n_0,\ldots,n_q)$, the number of objects is
\[
 \left(\frac{n!}{n_0!\cdots n_q!}\right)^2n_0!
 =\frac{(n!)^2}{n_0!\prod_{i=1}^q(n_i!)^2},
\]
which proves the coefficient interpretation.

\begin{corollary}\label{cor:family}
For every fixed positive integer $q$, the forward expansion is
\begin{align}\label{eq:familyforward}
 A_{q,n}\sim{}&\sqrt{2\pi}(4\pi)^{-q/2}
 n^{n+1/2-q/4}\e^{-n+2q\sqrt n-q^2/2}
 \sum_{j\ge0}d_{q,j}n^{-j/2},\qquad d_{q,0}=1.
\end{align}
Its late coefficients satisfy
\begin{equation}\label{eq:familylate}
 d_{q,j}\sim\frac{q\e^{2q-2}}\pi\frac{\Gamma(j)}{4^j}.
\end{equation}
All constants and remainder thresholds may depend on the fixed $q$.
\end{corollary}
\begin{proof}
Let $H_q=B^q$. Repeated endpoint convolution, using the bound from Section 3, gives $[t^j]H_q\sim q\Gamma(j)/(\pi4^j)$. The same forward saddle yields
\begin{align*}
 D_q(t)=\e^{q^2/2}S(t)\E\biggl[\e^{qu}(1+tu)^{-1-q/4}
 \e^{R_1(t,u)+R_{q,2}(t,u)}
 H_q\!\left(\frac{t}{\sqrt{1+tu}}\right)\biggr],\\
 R_{q,2}(t,u)=\sum_{m\ge1}2q\binom{1/2}{m+1}u^{m+1}t^m.
\end{align*}
The Bessel prefactor is $(4\pi)^{-q/2}x^{-q/4}$ and the linear Gaussian expectation is $\E\e^{qu}=\e^{-q^2/2}$, giving \eqref{eq:familyforward}. The exact outer-arc loss is still of order $n^{1/5}$, and multiplying the square-root phase by fixed $q$ does not compete with it.

At lag $\ell$, the amplitude rising factorial is $(k/2+1+q/4)_a$, bounded by $(C_q j)^a$. The proof of \eqref{eq:majorant}, including the $k=0$ endpoint, gives the summable majorant $C_q^\ell\ell^{-\ell/2}$. At fixed lag the normalized limit is $\E[\e^{qu}(-2u)^\ell/\ell!]$. Summing gives $\E\e^{(q-2)u}=\e^{-(q-2)^2/2}$; multiplication by $\e^{q^2/2}$ yields $\e^{2q-2}$. Stirling multiplication leaves the constant unchanged.
\end{proof}
For $q=2$ this is the main counting sequence; for $q=1$, $A_{1,n}=n!L_n(-1)$ follows from the finite Laguerre polynomial sum. No uniform estimate as $q$ grows, all-orders family late expansion, or separate family priority claim is made here.

\section{Reproducible checks and their interpretation}
The accompanying package contains the TeX source, generated results, and self-contained Python scripts. The mathematical argument supplies the asymptotic remainders; the numerical data illustrate them and test independent coefficient calculations. Finite numerical agreement is not used as a proof of any limiting statement or a certified error bound.

The replay runs in about seven seconds in the preparation environment using Python 3.12 and mpmath 1.3.0. It carries out the following tasks.
\begin{enumerate}
\item Construct $d_0,\ldots,d_{12}$ by exact rational Gaussian moments and compare them with a separately transformed recurrence. Check \eqref{eq:dsmall} exactly.
\item Compute $c_0,\ldots,c_8$ from the finite Bernoulli generator and compare every term with the independent untilted Gaussian transform through degree $t^{16}$ in \eqref{eq:cgaussian}. Odd Gaussian coefficients vanish exactly.
\item Recompute the late coefficients through index 100 at 200 and 300 working decimal digits, compare their stable digits, and evaluate the normalized leading ratios and four-term residuals.
\item Generate integer counts by the exact recurrence
\begin{align}\label{eq:arecurrence}
 n a_n={}&(3n^2+n-1)a_{n-1}-(n-1)^2(3n+1)a_{n-2}\nonumber\\
 &+(n-2)^2(n-1)^2a_{n-3}\qquad(n\ge3),
\end{align}
with $a_0=1,a_1=3,a_2=16$, compare small cases with \eqref{eq:an}, and evaluate the inverse formulas.
\end{enumerate}
The recurrence also follows by eliminating derivatives from the Bessel differential equation for the entire generating function in \eqref{eq:an}; the package checks its finite coefficients independently rather than treating a numerical fit as an exact identity.

\begin{table}[ht]
\centering
\begin{tabular}{rcc}
\toprule
$j$ & $d_j/[(2\e^2/\pi)4^{-j}\Gamma(j)]$ & four-term residual\\
\midrule
20 & $0.98814847357845965$ & $-3.06203\times10^{-6}$\\
40 & $0.99720458884688425$ & $-3.36102\times10^{-11}$\\
60 & $0.99878530659374644$ & $-1.17552\times10^{-12}$\\
80 & $0.99932469727066581$ & $-1.05220\times10^{-13}$\\
100& $0.99957088850853692$ & $-1.63631\times10^{-14}$\\
\bottomrule
\end{tabular}
\caption{High-precision numerical checks, not interval-certified bounds. The residual subtracts the four displayed terms in \eqref{eq:normalized}. Fixed-order asymptotics do not promise monotone improvement at each small index.}\label{tab:late}
\end{table}

\begin{table}[ht]
\centering
\begin{tabular}{rcc}
\toprule
$n$ & $r_A(a_n)-n$ & $(r_A(a_n)-n)N\log N$\\
\midrule
20 & $4.87410\times10^{-3}$ & $0.410219$\\
50 & $5.60735\times10^{-4}$ & $0.129746$\\
100 & $1.93015\times10^{-6}$ & $0.000983211$\\
200 & $-7.60224\times10^{-5}$ & $-0.0856808$\\
500 & $-5.03700\times10^{-5}$ & $-0.161750$\\
1000 & $-2.85666\times10^{-5}$ & $-0.201469$\\
\bottomrule
\end{tabular}
\caption{Smooth inverse checks against exact integer counts. The scale is consistent with \eqref{eq:raerror}; these rows do not determine the unknown enclosure constant or justify rounding at arbitrary targets.}\label{tab:inverse}
\end{table}

The two precision runs differ by at most $1.04\times10^{-196}$ under the metric $|d_j^{(200)}-d_j^{(300)}|/\max(1,|d_j^{(300)}|)$. The package also compares the late Lambert approximation with a numerically solved four-term Gamma model. At $j=100$ their inverse errors against the computed coefficient target are approximately $-1.91\times10^{-4}$ and $-5.09\times10^{-15}$, respectively. The model root is not an exact interpolation of the discrete sequence.

The package's README gives one-command replay and PDF rebuild instructions. The manifest checks the distributed files, and the replay compares regenerated result files. All paths are relative. The exact arithmetic outputs are exact; stability under higher precision is evidence about the floating-point calculations, not an interval certificate.

\section{Sources and boundaries of the conclusion}
The input definitions and conjecture were read from accessible indexed OEIS entry content on 2 October 2026. Direct entry, sequence-text and b-file retrieval attempts did not produce a fresh downloadable source snapshot: direct requests returned access errors and b-file requests failed. The exact finite checks here therefore validate the internally derived rational coefficients and displayed entry examples, not raw OEIS b-files. The package includes source links and a short retrieval account, not third-party full texts.

A bounded overlap search inspected 185 individually retrieved complete TeX or Markdown texts from the ProveIt repository at commit
\texttt{4b874cea0012c51a6841ad9c58f6e20fa57c71da}. It included selected generating-function and asymptotic texts, thematic factorial and Borel material, and the canonical transseries volumes. No proof of this target conjecture was located in that inspected set. Relevant Bessel and Laguerre passages concerned different sequences or problems. This is a bounded content comparison, not a claim that every repository file, historical version or binary was searched, and not a worldwide priority claim.

The analytic ingredients are standard. Bessel large-argument expansions and their product identity are classical; Stirling, beta and reciprocal-Gamma integrals provide the remaining special-function tools. The product and composition laws used here are Borinsky's established results. The work specific to this sequence is the connection to its actual forward saddle, the uniform Gaussian transfer estimates, and the resulting explicit identification of its late corrections.

Several questions remain beyond the established scope:
\begin{enumerate}\setlength{\itemsep}{0pt}\setlength{\parsep}{0pt}
\item Prove or disprove positivity for every $j>5$, with a genuinely effective finite check and eventual bound.
\item Turn the fixed-order remainders into explicit usable constants and cutoffs, including a certified integer inverse algorithm at a specified finite target.
\item Determine the arithmetic growth of the reduced numerator and denominator separately; quotient growth alone does not determine either.
\item Analyze optimal truncation, exponentially small errors, or global Borel continuation. The present all-orders fixed-order theorem does not assert any of these stronger facts.
\end{enumerate}

\begingroup\small
\begin{thebibliography}{9}
\setlength{\itemsep}{2pt}
\bibitem{oeis336} OEIS Foundation, \emph{A336293}, \url{https://oeis.org/A336293}. Indexed entry consulted 2 October 2026.
\bibitem{oeis395976} OEIS Foundation, \emph{A395976}, \url{https://oeis.org/A395976}. Rational numerators and conjecture attributed in the indexed entry to V. Kotesovec, 21 May 2026.
\bibitem{oeis395977} OEIS Foundation, \emph{A395977}, \url{https://oeis.org/A395977}. Corresponding rational denominators.
\bibitem{dlmfbessel} NIST Digital Library of Mathematical Functions, \S10.40, especially 10.40.1, 10.40.6 and \mbox{\S10.40(iii)}, \url{https://dlmf.nist.gov/10.40}.
\bibitem{dlmfgamma} NIST Digital Library of Mathematical Functions, 5.9.2, reciprocal-Gamma Hankel integral, \url{https://dlmf.nist.gov/5.9.E2}.
\bibitem{dlmfstirling} NIST Digital Library of Mathematical Functions, \S5.11, especially 5.11.8, shifted logarithmic Gamma expansion, \url{https://dlmf.nist.gov/5.11.E8}.
\bibitem{borinsky} M. Borinsky, \emph{Generating asymptotics for factorially divergent sequences}, Electron. J. Combin. \textbf{25}(4) (2018), Paper 4.1. Proposition 22, Theorem 35 and Lemma 37; version 4, \url{https://arxiv.org/html/1603.01236v4}.
\end{thebibliography}
\endgroup
\end{document}
